\DocumentMetadata{
  pdfversion=2.0,
  pdfstandard=ua-2,
  lang=en-US,
  tagging=on,
  tagging-setup={math/setup=mathml-SE,tabsorder=structure}
}
\documentclass[11pt,letterpaper]{article}
\usepackage[margin=0.68in,includefoot]{geometry}
\usepackage{newcomputermodern}
\usepackage{amsmath,amscd,mathtools}
\usepackage{tikz,adjustbox}
\providecommand{\square}{\mdlgwhtsquare}
\usepackage[hidelinks]{hyperref}
\hypersetup{
  pdftitle={Numerical Analysis PhD Exam, September 2005},
  pdfauthor={University of Florida Department of Mathematics},
  pdfsubject={Graduate mathematics examination},
  pdfdisplaydoctitle=true,
  bookmarksopen=true
}
\setcounter{secnumdepth}{0}
\setlength{\parindent}{0pt}
\setlength{\parskip}{0.28em}
\setlength{\emergencystretch}{3em}
\setlength{\leftmargini}{2.15em}
\setlength{\leftmarginii}{2.65em}
\setlength{\leftmarginiii}{3em}
\renewcommand{\labelenumi}{\textbf{\theenumi.}}
\pagestyle{empty}
\raggedbottom
\allowdisplaybreaks
\newenvironment{examproblems}[1]{%
  \begin{enumerate}%
  \setcounter{enumi}{#1}%
  \setlength{\topsep}{0.35em}%
  \setlength{\partopsep}{0pt}%
  \setlength{\itemsep}{0.48em}%
  \setlength{\parsep}{0.18em}%
}{\end{enumerate}}
\newenvironment{examsubparts}{%
  \begin{enumerate}%
  \setlength{\topsep}{0.18em}%
  \setlength{\partopsep}{0pt}%
  \setlength{\itemsep}{0.16em}%
  \setlength{\parsep}{0.1em}%
}{\end{enumerate}}
\newenvironment{examinstructions}{%
  \par\begingroup\itshape
}{%
  \par\endgroup\vspace{0.18em}
}
\makeatletter
\renewcommand\section{\@startsection{section}{1}{0pt}%
  {-1.25ex plus -.25ex minus -.1ex}%
  {0.55ex plus .1ex}%
  {\normalfont\large\bfseries}}
\renewcommand{\@maketitle}{%
  \newpage\null\vskip -1em%
  {\footnotesize\bfseries
    \href{https://www.ufl.edu/}{University of Florida}, \href{https://math.ufl.edu/}{Department of Mathematics}\par}%
  \vskip -0.5em%
  \tagstructbegin{tag=Title}%
    {\LARGE\bfseries\href{https://gma.math.ufl.edu/exams/numerical-analysis-phd/}{Numerical Analysis PhD Exam}, September 2005\par}%
  \tagstructend%
  \vskip 0.25em%
  \tagmcbegin{artifact}%
    \hrule height 1pt%
  \tagmcend%
  \vskip 1em%
}
\makeatother
\title{Numerical Analysis PhD Exam, September 2005}
\author{}
\date{}
\begin{document}
\maketitle
\thispagestyle{empty}
\begin{examinstructions}
Do any eight of the ten problems below.
\end{examinstructions}
\begin{examproblems}{0}
\item
For solving \(dy/dt=f(t,y)\), consider a one-step method of the form 
\[
y_1=y_0+h\Phi(t_0,y_0,h).
\]
 For any \(t_0\), it computes an approximation \(y_1\) to \(y(t_0+h)\), given \(y(t_0)=y_0\). Assuming enough smoothness on~\(f\), study the local truncation error of Runge–Kutta methods for which 
\[
\Phi(t_0,y_0,h)=\alpha_1k_1+\alpha_2k_2,
\]
 where 
\[
k_1=f(t_0,y_0),
\]
 and 
\[
k_2=f(t_0+\mu h,y_0+\mu hk_1).
\]
 Under what conditions on the constants \(\alpha_1,\alpha_2\), and~\(\mu\) do we get second order methods?
\item
Prove that any polynomial~\(q\) of degree at most \(n-1\) satisfies 
\[
\sum_{i=0}^n q(x_i)\prod_{\substack{j=0\\j\ne i}}^n\frac{1}{x_i-x_j}=0,
\]
 for any distinct \(n+1\) numbers \(x_i\), \(i=0,1,\ldots,n\). (Suggestion: Identify the left hand side as a divided difference.)
\item
Let \(p_2(x)\) be the quadratic polynomial that interpolates \(f(x)\) at nodes \(x=0,h\), and \(2h\).
\begin{examsubparts}
\item[\textnormal{(a)}]
Approximating 
\[
Q=\int_0^{3h}f(x)\,dx
\]
 by 
\[
Q_h=\int_0^{3h}p_2(x)\,dx,
\]
 derive a quadrature rule: Explicitly calculate \(w_1,w_2,w_3\) so that 
\[
Q_h=w_1f(0)+w_2f(h)+w_3f(2h).
\]
\item[\textnormal{(b)}]
Assuming that \(f(x)\) is four times continuously differentiable, show that 
\[
Q-Q_h=\frac{3}{8}h^4f'''(0)+O(h^5).
\]
\end{examsubparts}
\item
Denote the complete (or clamped) and natural cubic splines interpolating \(f\in C^2[a,b]\) at knots \(t_0<t_1<\cdots<t_n\) by \(s_c(t)\) and \(s_n(t)\), respectively.
\begin{examsubparts}
\item[\textnormal{(a)}]
Prove that 
\[
\int_a^b|s_c''(x)|^2\,dx\leq\int_a^b|f''(x)|^2\,dx.
\]
 State an analogous inequality for \(s_n\).
\item[\textnormal{(b)}]
Which of the two interpolating splines has smaller “energy”, i.e., which of \(\int_a^b|s_c''(x)|^2\,dx\) and \(\int_a^b|s_n''(x)|^2\,dx\) is smaller?
\end{examsubparts}
\item
Let \(p>0\) and 
\[
x=\sqrt{p+\sqrt{p+\sqrt{p+\cdots}}},
\]
 where all the square roots are positive. Let \(F(y)=\sqrt{p+y}\). (How is the fixed point of~\(F\) related to~\(x\)?)
\begin{examsubparts}
\item[\textnormal{(a)}]
Consider the fixed point iteration \(x_{n+1}=F(x_n)\). Prove that if the initial iterate \(x_0\) satisfies \(x_0+p>0\) then all iterates remain on the same side of~\(x\) as \(x_0\).
\item[\textnormal{(b)}]
Prove that the fixed point iteration converges for all choices of initial guesses greater than \(-p+1/4\).
\item[\textnormal{(c)}]
Whenever the fixed point iteration converges, what is its order of convergence?
\end{examsubparts}
\item
This problem has three parts.
\begin{examsubparts}
\item[\textnormal{(a)}]
Consider the matrix \(A=uv^*\) where~\(u\) and \(v\in\mathbb{C}^n\). Under what condition on~\(u\) and~\(v\) is~\(A\) a projector?
\item[\textnormal{(b)}]
Show that the Householder matrix \(H=I-2ww^*\) where \(\lVert w\rVert=1\) is unitary.
\item[\textnormal{(c)}]
Given a vector \(x\in\mathbb{C}^n\) and an integer~\(k\) with \(1<k<n\), derive a formula for a Householder matrix with the property that \((Hx)_i=0\) for \(i>k\) and \((Hx)_i=x_i\) for \(i<k\). Be sure to choose the signs so that the formula is numerically stable.
\end{examsubparts}
\item
Consider the conjugate gradient algorithm for solving \(Ax=b\) for a symmetric positive definite \(A\in\mathbb{R}^{m\times m}\).
\begin{examsubparts}
\item[\textnormal{(a)}]
State a relationship between the~\(n\)-th iterate of the algorithm and the best approximation to~\(x\) from a Krylov space (be sure to specify the norm).
\item[\textnormal{(b)}]
Prove that if~\(A\) has only \(n\leq m\) distinct eigenvalues then the iteration (in the absence of round off errors) converges to the exact solution~\(x\) in at most~\(n\) steps no matter what the initial iterate is. (You may use 7a.)
\end{examsubparts}
\item
Let~\(P\) and~\(Q\) be two \(m\times m\) orthogonal projectors.
\begin{examsubparts}
\item[\textnormal{(a)}]
Prove that \(\lVert P-Q\rVert_2\leq 1\).
\item[\textnormal{(b)}]
Prove that if \(\lVert P-Q\rVert_2<1\) then the ranges of~\(P\) and~\(Q\) have equal dimensions.
\end{examsubparts}
\item
Let~\(P\) and~\(Q\) be Hermitian positive definite matrices. Prove that 
\[
x^*Px\leq x^*Qx,\qquad\text{for all }x\in\mathbb{C}^n,
\]
 if and only if 
\[
x^*Q^{-1}x\leq x^*P^{-1}x,\qquad\text{for all }x\in\mathbb{C}^n.
\]
\item
State the Rayleigh quotient iteration. Apply it to the \(2\times2\) real matrix 
\[
A=\begin{bmatrix}\lambda_1&0\\0&\lambda_2\end{bmatrix},
\]
 where \(\lambda_1\ne\lambda_2\). Find the subset \(S\subseteq\mathbb{R}^2\) having the property that the iteration applied to this matrix with initial guesses in~\(S\) do not converge. Is~\(S\) a set of measure zero?
\end{examproblems}
\end{document}
